% PAGES: 66  (printed 50-50)
% Opens mid-"examplex" environment: closes the $4\times4$-matrix example
% begun in sec02_c2_3.tex (page 49). Do not add a new \begin{examplex}
% here -- the \end{examplex} below matches the \begin{examplex} in
% sec02_c2_3.tex.
% This file stops at the end of printed page 50, right before the
% "for k = i:n ... end" pseudocode box, which is the first content of
% page 51 (out of range for this agent).

\[
U(1,2:n) \;=\; U(1,2:4) \;=\; \big(U(1,2)\ \ U(1,3)\ \ U(1,4)\big);
\]
\[
U(2:n,2:n) \;=\; U(2:4,2:4) \;=\;
\begin{pmatrix} U(2,2) & U(2,3) & U(2,4)\\ 0 & U(3,3) & U(3,4)\\ 0 & 0 & U(4,4)\end{pmatrix},
\]
and therefore (2.28--2.29) is written as follows if $A$, $L$, and $U$
are $4\times 4$ matrices:
\begin{equation}
\begin{pmatrix} 1 & 0 \\ L(2:4,1) & L(2:4,2:4)\end{pmatrix}
\begin{pmatrix} U(1,1) & U(1,2:4) \\ 0 & U(2:4,2:4)\end{pmatrix}
\end{equation}
\begin{equation}
\;=\; \begin{pmatrix} A(1,1) & A(1,2:4) \\ A(2:4,1) & A(2:4,2:4)\end{pmatrix}.
\end{equation}

From (2.33--2.34), we find by block matrix multiplication the following
explicit form of (2.30) for $4\times 4$ matrices:
\[
L(2:4,1)U(1,2:4) \;+\; L(2:4,2:4)U(2:4,2:4) \;=\; A(2:4,2:4),
\]
and therefore
\begin{equation}
L(2:4,2:4)U(2:4,2:4) \;=\; A(2:4,2:4) \;-\; L(2:4,1)U(1,2:4).
\end{equation}
\end{examplex}

The LU decomposition without pivoting algorithm can be implemented
recursively as detailed below; see also the pseudocode from Table 2.5.

Let $A$ be a matrix with LU decomposition without pivoting, and let $L$
and $U$ be the LU factors of $A$.

\begin{itemize}
\item[$\bullet$] Compute the first row of $U$ and the first column of $L$:
\end{itemize}

\begin{center}
\fbox{%
\begin{minipage}{0.5\textwidth}
for $k = 1:n$\\
\hspace*{1.5em}$U(1,k) = A(1,k)$\\
\hspace*{1.5em}$L(k,1) = A(k,1)/U(1,1)$\\
end
\end{minipage}%
}
\end{center}

Note that, in the ``for'' loop above, we compute $L(1,1) = 1$, rather
than requiring it.

\begin{itemize}
\item[$\bullet$] Update the $(n-1)\times(n-1)$ lower right part of $A$ as follows:
\end{itemize}
\begin{equation}
A(2:n,2:n) \;=\; A(2:n,2:n) \;-\; L(2:n,1)U(1,2:n),
\end{equation}
see (2.32), which can be written entry by entry as follows:

\begin{center}
\fbox{%
\begin{minipage}{0.55\textwidth}
for $j = 2:n$\\
\hspace*{1.5em}for $k = 2:n$\\
\hspace*{3em}$A(j,k) = A(j,k) - L(j,1)U(1,k)$\\
\hspace*{1.5em}end\\
end
\end{minipage}%
}
\end{center}

Every row of $U$ and column of $L$ are thereafter computed recursively
from the latest updated part of the matrix $A$; see also the $4\times4$
example below. For example, to compute the $i$-th row of $U$ and the
$i$-th column of $L$, we do the following:

\begin{itemize}
\item[$\bullet$] Compute the $i$-th row of $U$ and the $i$-th column of $L$:
\end{itemize}
